%% Escalate or Answer -- living paper for CS 567 checkpoints (ACM CHI format).
%% Checkpoint 2 (methodology): single-column review layout, as the course requires.
\documentclass[manuscript,review]{acmart}

\settopmatter{printacmref=false, printccs=false, printfolios=true}
\setcopyright{none}
\acmConference{CS 567: 3D User Interfaces and Human-Centered Spatial AI}{Fall 2026}{Colorado State University}

\usepackage{xcolor}
\usepackage{graphicx}
\usepackage{booktabs}
\setlength{\emergencystretch}{1em}

%% Visible marker for claims the authors must still check against the cited paper.
%% Replace the definition with \newcommand{\verify}{} before a submission that should not show them.
\newcommand{\verify}{\textsuperscript{\textcolor{red}{[verify]}}}
%% Placeholder for sections that are not yet written.
\newcommand{\todo}[1]{{\itshape\color{black!60}[#1]}}

\begin{document}

%% ---------------------------------------------------------------- title
\title{How Should an Agentic RAG Assistant Hand Back a Decision That Depends on the User?}

\author{Coen Petto}
\email{coenp@colostate.edu}
\affiliation{%
  \institution{Colorado State University}
  \city{Fort Collins}
  \state{Colorado}
  \country{USA}
}
\author{Cameron Suess}
\email{camsuess@colostate.edu}
\affiliation{%
  \institution{Colorado State University}
  \city{Fort Collins}
  \state{Colorado}
  \country{USA}
}

\renewcommand{\shortauthors}{Petto and Suess}

%% ---------------------------------------------------------------- abstract
\begin{abstract}
Retrieval-augmented generation (RAG) assistants over internal documentation
answer fluently even when the right answer depends on who is asking. Agentic
RAG architectures prescribe a remedy: when the agent cannot settle a choice
from the documents alone, it hands the decision back to the person. How that
hand-back should look is untested where the agent's recommendation can be
wrong. We ran a $3\times2$ within-subjects study with [$N$] university
students who were new to their department's GPU cluster. Each participant
held a spec sheet for a course project and, on each of nine tasks, sent a
scripted query, watched the assistant's retrieval trace stream over a frozen
documentation snapshot, and accepted or rejected the single answer it filled
in. Every task was a fork the documents could not settle, and the deciding
fact was on the sheet, which the assistant never saw. We varied only the
form of the hand-back at the fork (modal stop, inline warning, or
provenance-only citation), holding content, provenance and undo constant,
and measured which sources were opened and for how long, whether
participants kept or overrode the recommendation when it was right and
wrong, decision time, and workload. \todo{Results and implications are
added after data collection.}
\end{abstract}

\keywords{agentic AI; retrieval-augmented generation; interruption; appropriate reliance; human-in-the-loop; proactive assistants}

\maketitle

%% ---------------------------------------------------------------- body
\section{Introduction}
\label{sec:intro}

Retrieval-augmented generation (RAG) is the usual way to put a large
language model on top of an organization's own documents, and internal
documentation is full of answers that are conditional. On the Colorado
State University (CSU) Computer Science department's Falcon cluster, the
quality-of-service (QoS) level a job is submitted under depends on how long
the job will run: the documentation lists a \texttt{gpu\_short} QoS with a
24-hour limit and a \texttt{gpu\_medium} QoS with a 3-day limit, and which
one applies depends on the job. How to connect depends on whether you are
on campus; where a job should write depends on whether the data matters.
People ask short questions that leave the deciding fact out (``Which QoS
should I use to submit this job?''), and a RAG assistant composes a
confident answer from whichever passage ranked higher. Nothing on screen
says that the other passage applies to someone else, and a reader who is
new to the system has no way to tell.

The reliance literature gives reason to expect that the reader will not
catch it. Bansal et al.~\cite{bansal2021} found that explanations raise
acceptance of AI recommendations whether or not they are right.
Bu\c{c}inca et al.~\cite{bucinca2021} showed that cognitive forcing
functions reduce overreliance but are the designs people like least.
Vasconcelos et al.~\cite{vasconcelos2023} modeled verification as a
cost--benefit choice: people check the AI only when checking is cheap
relative to the task. The practical question is therefore not whether to
show sources but how to hand the decision back at the moment it is needed.

Agentic RAG architectures describe such a moment. Mishra et
al.~\cite{mishra2026} catalogue a human-in-the-loop pattern in which the
agent pauses execution to request disambiguation when it cannot resolve a
choice itself. S\'anchez-Adame and Mendoza~\cite{sanchezadame2026} supply
a contract, Minimal Proactivity Policies (MiPP), for how a system-initiated
pause should be expressed, and tested three initiation forms in a
within-subjects study of $n{=}48$, but on a system that was never wrong,
so checking the source could not change any outcome. This paper asks what
neither line has answered:

\begin{quote}
\textbf{RQ.} When an agentic RAG assistant finds mid-task that the right
answer depends on a fact about the user that the question did not state,
and hands the choice back, how does the form of that hand-back (modal
stop, inline warning, or provenance-only citation) affect verification,
appropriate reliance, decision time and workload?
\end{quote}

We hold the agent constant and vary only the human-facing form. Each task
is one scripted assistant session over a frozen snapshot of the Falcon
documentation; the participant watches the retrieval trace, may open
either of the two sources it cites, and accepts or rejects the one answer
it fills in. Every task is a \emph{context-dependent fork}: two passages
give different answers, both correct for someone, and which applies
depends on the participant's situation, written on a spec sheet the
assistant was never given. Such a fork is the case the human-in-the-loop
pattern is written for: the agent cannot settle the choice from the
documents and the person can, so pausing to ask is the pattern's intended
use, not a failure~\cite{mishra2026,horvitz1999}. The recommendation is wrong for this
participant on about half of the tasks, so verifying has a payoff and keeping
or overriding it can be scored as appropriate or inappropriate
reliance~\cite{schemmer2023}.

This paper makes three contributions. (1) A comparison of three
escalation forms at fixed decision points on an agent whose recommendation
can be wrong, with verification and appropriate reliance as outcomes. (2)
Predictions, derived from the interruption and reliance literatures, of
how the forms trade verification against decision time and workload, and
their test. (3) A reusable protocol (frozen snapshot, task records,
contamination filter, matched task sets, event logging) that other groups
can apply to their own documentation.

\section{Related Work}
\label{sec:related}

\subsection{Interruption as a design variable}
McFarlane and Latorella~\cite{mcfarlane2002scope} made human interruption
a first-class problem for interface design.
McFarlane~\cite{mcfarlane2002comparison} compared four coordination
methods (immediate, negotiated, mediated, scheduled) in a dual-task
experiment with 36 participants and concluded that ``negotiation support
is the best overall solution except where small differences in the
timeliness of handling interruptions is critical and then immediate is
best.''
Horvitz~\cite{horvitz1999} framed mixed initiative decision-theoretically:
initiate only when the expected utility of acting exceeds the cost of
interrupting, and use dialogue to resolve uncertainty. Read against this
work, and as our own interpretation, a modal stop is closest to
McFarlane's immediate method, an inline warning to a negotiated one, and a
provenance-only answer does not interrupt at all. \verify{} In these studies the
interruption carried no fallible content; what they leave open is whether
the coordination method also changes the accuracy of the decision the
interruption asks the person to make.

\subsection{Reliance on AI advice}
Work at CHI, CSCW and IUI studies what people do with an AI recommendation
that may be wrong. Bansal et al.~\cite{bansal2021} found no complementary
gain from explanations, which raised acceptance regardless of correctness.
Bu\c{c}inca et al.~\cite{bucinca2021} showed that forcing engagement before
the AI's answer is accepted reduces overreliance at a subjective cost.
Vasconcelos et al.~\cite{vasconcelos2023} showed across five studies that
explanations help only when they lower the cost of verifying relative to
the task. Schemmer et al.~\cite{schemmer2023} define appropriateness of
reliance in two dimensions: following the AI when it is right and
overriding it when it is wrong. M\"ohle et al.~\cite{mohle2026} found in a
judge--advisor study that a stated accuracy level shifted self-reported
trust more when the reasoning was convincing, while actual accuracy was
never varied. We take from this that self-report is a weak proxy for
behavior, which is why our primary outcomes are logged. All of these are
single-shot decisions in which the AI proposes and the human disposes and none
has a system-initiated pause mid-task, and none varies its form.

\subsection{MiPP: a contract for proactive initiation}
MiPP~\cite{sanchezadame2026} separates decision authority (whether the
system may initiate) from interaction authority (how the initiation is
expressed). Its constraints include C1, a brief cue followed by a single
question or choice; C2, a one-step escape hatch; and C3, an explicit
provenance cue naming the source that justified the initiation. C5 says
that when inputs are missing or ambiguous the system should ``defer,
downgrade to an indirect channel, or remain silent.'' MiPP's $n{=}48$
study compared a reactive baseline with three policy families (direct
alert, confirmatory prompt, indirect entry) on productivity scenarios; the
quieter families reduced perceived intrusion, and removing the
confirmatory cue worsened intrusion, timing, trust and willingness. Three
features of that study define the gap we address. Each condition bundled a
family with its own trigger and scenario, so forms were not compared at the
same decision points. Every initiation was correct by construction, so
verification had no consequence and the families differed on it by at most
six points. Its verification measure is anchored to an initiation, so
it is undefined when nothing is announced. We anchor verification to the
task. MiPP calls its contract domain-agnostic and names per-scenario
designs with real workloads as future work.

C5 also creates a tension that our design tests. In MiPP, silence means
non-initiation: when the inputs that describe the user's situation are
missing, the system does not initiate. Our agent is answering a query it
was asked, so the move it can decline to initiate is the hand-back at the
fork, and the fact it lacks is exactly the kind C5 names, a fact about the
user. The provenance-only form is therefore non-initiation in C5's sense:
the assistant answers and cites but does not raise the fork. In MiPP's
setting that default is safe because nothing is lost by staying quiet;
here it means quietly applying a recommendation that is wrong for this
user.

\subsection{Agentic RAG and RAG transparency}
Mishra et al.~\cite{mishra2026} systematize agentic RAG and document a
human-in-the-loop pattern that ``models human oversight as a callable API
within the action space,'' pausing execution ``to request disambiguation
or supervision'' when uncertainty exceeds a threshold. They note that
existing implementations validate at predefined bottlenecks rather than
triggering on internal ambiguity, and \emph{propose} escalation precision
and recall and detection of conflicting retrieved contexts as criteria
for future systems; the paper is a survey, available only as an arXiv
preprint, and evaluates nothing. Ravi and Sindhgatta~\cite{ravi2025}
evaluated a RAG assistant for equity research with 50 financial
professionals and report that source attribution and highlighting
substantially improved self-reported trust and understanding, whereas
confidence scores alone did not.
Their measures are self-report, the system is not agentic, and nothing
interrupts mid-task; a cited answer with named sources is the current
practice that our provenance-only condition reproduces. Kirchenbauer et
al.~\cite{kirchenbauer2026} validated FictionalQA items by prompting the
generator blind and informed and discarding items answerable blind; the
check reduces, not removes, the pretraining confound, and we reuse it as a
filter on real Falcon facts.

\subsection{Pre-generated agent behavior in reliance studies}
Studies of reliance on AI advice typically fix the AI's advice in advance
and show every participant the same items, so that the correctness of each
recommendation is controlled rather than left to a live
model~\cite{bucinca2021,vasconcelos2023}; the practice has the same aim
as Wizard-of-Oz studies of intelligent interfaces, in which system
behavior is controlled so that the interface can be studied on its
own~\cite{dahlback1993}. We follow that paradigm: the
agent's behavior is generated once by the pipeline described in
Section~\ref{sec:method} and replayed identically to every participant.

\subsection{The gap}
Interruption research tells us how coordination method changes disruption
when the interruption is not fallible. Reliance research tells us how
forcing engagement changes overreliance in single-shot decisions. MiPP
tells us how an initiation should be expressed, on a system that is never
wrong. No study we are aware of has fixed the decision points of a fallible
agent, varied only the form in which it hands a decision back mid-task,
and measured whether people verify and rely appropriately, and at what
cost.

%% ---------------------------------------------------------------- method
\section{Method}
\label{sec:method}

%% Checkpoint 2 focus. Written in past tense, as if the study has been run.
%% Subsection order follows MacKenzie (2024), Ch. 5 and 8.2.5.

\subsection{Participants}
\label{sec:participants}
Twelve Colorado State University students who reported no prior use of
the department's Falcon cluster took part ([$n$] undergraduate, [$n$]
graduate; [$n$] women, [$n$] men; courses of study were recorded at the
session). They were recruited in person by word of mouth, were not
compensated, and gave consent on the course's class-project
data-collection form. Three pilot participants run before data collection
are not included.

\subsection{Apparatus}
\label{sec:apparatus}
\paragraph{Corpus.} The corpus was a frozen snapshot of the CSU Computer
Science department's systems documentation (sna.cs.colostate.edu), crawled
on 3 October 2026: 77 pages (23,965 words), stored as raw HTML and
Markdown with per-page hashes, split at headings into 226 chunks and
embedded with \texttt{nomic-embed-text}.

\paragraph{Drafting pipeline.} The assistant's retrieval and escalation
decisions were produced before the study by a local pipeline
(\texttt{qwen2.5:7b-instruct} through Ollama 0.35.1; temperature 0, seed
567) and frozen as task records, so no model ran during a session. \verify{} Three
elements of each task came from the pipeline: the retrieved passages and
their ranking (an embedding search over the snapshot for the scripted
query, top four chunks); the decision that the question was a fork, made
by the fork check, which fired when the cosine margin between the top two
chunks was at most 0.10 and an LLM judge found that the passages gave
different correct answers depending on a fact the query did not state; and
the recommendation, which was always the option whose passage ranked
higher. \verify{} The two options, the highlighted passage for each, the
spec-sheet lines and the answer keys were written by us from the snapshot.
The pipeline then had to find the same fork on its own: shown the
retrieved passages, the judge named the options it saw, and a fork was
used only when those options matched the two we had authored, which held
for 8 of 25 query phrasings over 15 candidate forks. The seven-line trace
shown to participants was a fixed template; the pipeline supplied what
filled it: which page was Source A and which Source B (its retrieval
ranking), that line 6 was a fork (its fork check), and which option line 6
recommended (the higher-ranked option). The wording of the lines
themselves (``Reading the question\ldots'', ``Found Source A: \ldots'',
``Source A suggests [option].'') was fixed, not written by the model, so
that every task presented the same structure and the form of the fork
line was the only variable.

\paragraph{Interface.} \verify{} The study ran as a single-page web application in a
desktop browser: the block's spec sheet, the task question, one answer
field and Accept and Reject buttons on the left; the scripted query, a Send
button and the assistant's seven-line trace on the right, in which Source A
and Source B appeared as chips that opened the frozen page with the cited
lines highlighted. Trace lines arrived one at a time at a fixed delay of
3--5 s stored in the task record, and only line 6, the fork line, was
rendered differently across forms. Instrument screens between blocks and
at the end of the session collected the questionnaires, and the
application assigned each participant's block order and task rotation
from the participant number.

\paragraph{Logging.} Every interaction (Send, each line onset, fork onset,
modal pick, each source open and close with its duration, Reject, picker
choice, Accept) was written to a local SQLite database under a participant
code, together with a per-task summary of the final value, accuracy,
following, verification, source pattern, dwell per source and decision
time.

\subsection{Tasks and Materials}
\label{sec:tasks}
The task pool held nine context-dependent forks, three per set (job
submission; access and remote connection; storage and software), plus two
practice tasks without a fork. A fork was two documented options whose
correct choice depended on one fact the scripted query omitted and the
spec sheet stated. Each fork existed in two twin versions with the same
query, trace, sources and recommendation; the twins differed only in the
sheet line and the task question, so that the recommendation was correct
under one and wrong under the other. Each participant met every fork
exactly once, in one of its two versions, one task from each set in every
block, with the assignment of forks to forms and of twins to participants
rotated so that every fork appeared under every form, and in each version,
equally often. The cover story was a single
course project on the cluster, and each block's spec sheet listed the
three deciding facts for that block's tasks, one line per task. The
recommendation was always the option whose passage ranked higher in
retrieval, and the key was the option the documentation gave for the
sheet line; correctness was never set by editing the agent.

\subsection{Design}
\label{sec:design}
The experiment was a $3\times2$ within-subjects design with factors
Escalation Form (S1 modal stop, S2 inline warning, S3 provenance-only) and
Recommendation Correctness (correct, wrong). Each form was administered in
one block of three tasks, one from each task set, so that task set was
crossed with form rather than tied to it; the twin assignment was
constrained so that every block contained at least one correct and one
wrong recommendation and every participant received four or five of each.
Block order was counterbalanced across the six permutations of the three
forms, two participants per order, and the fork-to-form rotation was
balanced within each order. With 12 participants this gave 108 measured
tasks, 18 per Form $\times$ Correctness cell.

\subsection{Procedure}
\label{sec:procedure}
\verify{} Sessions lasted about one hour: instructions stating that the
assistant was usually right but could be wrong and that any cited source
could be opened; two practice tasks, during which the experimenter had
participants open a source and use Reject and Accept once each; three
blocks of three tasks, one block per form, each block's three-line spec
sheet staying on screen; the instrument screens after each block; and a
closing ranking screen followed by a debrief. On each task the
participant pressed Send, watched the trace, could open either source,
and pressed Accept, or Reject followed by a choice from the two-option
picker; no feedback on correctness was given.

\subsection{Measures}
\label{sec:measures}
Per task we recorded verification (a source viewed for at least 3 s of
summed dwell before Accept), the source pattern under the same threshold
and the raw dwell time per source, following (final value equal to
the recommendation), accuracy (final value equal to the key) and decision
time (fork-line onset to Accept). After each block an instrument screen
collected raw NASA-TLX, a single perceived-intrusion item (``How
intrusive was the assistant's interruption?'', 1--7, the item used in
MiPP-Eval, kept for comparability with MiPP), and a manipulation check
asking which form the participant had just used (paused and asked, warned
but continued, or cited only), and at the end of the
session a screen collected a ranking of the three forms with a free-text
reason.

\subsection{Hypotheses and Analysis}
\label{sec:analysis}
We predicted that verification would be ordered S1 $>$ S2 $>$ S3 (H1),
because designs that force engagement before an answer can be accepted
raise engagement with the evidence~\cite{bucinca2021} and an immediate
interruption is handled most promptly~\cite{mcfarlane2002comparison};
that overriding a wrong recommendation would be higher under S1 and S2
than under S3 while following a correct one would not differ (H2, a Form
$\times$ Correctness interaction), because forcing functions reduce
overreliance~\cite{bucinca2021} whereas a cited answer alone raises
acceptance regardless of correctness~\cite{bansal2021}, scored as
relative self-reliance and relative AI reliance~\cite{schemmer2023}; that
decision time would be ordered S1 $>$ S2 $>$ S3 (H3a), because an
immediate interruption costs the interrupted task the
most~\cite{mcfarlane2002comparison} and checking a source is a time cost
people weigh~\cite{vasconcelos2023}; and that perceived workload would
follow the same order (H3b), because forcing functions are rated as more
effortful~\cite{bucinca2021} and quieter initiations as less
intrusive~\cite{sanchezadame2026}. Task-level outcomes (H1, H2, H3a) were
analyzed with logistic mixed models with fixed effects of form,
correctness, their interaction and block position, and random intercepts
for participant and task; decision time with a linear mixed model on log
time; and the per-block scales (H3b) with repeated-measures ANOVA.

%% ---------------------------------------------------------------- results
\section{Results}
\label{sec:results}
\todo{To be written after data collection.}

\subsection{Verification}
\subsection{Appropriate reliance}
\subsection{Decision time and workload}
\subsection{Subjective ranking and reasons}

%% ---------------------------------------------------------------- discussion
\section{Discussion}
\label{sec:discussion}
\todo{To be written after data collection.}

\subsection{Implications for proactive assistants}
\subsection{Limitations}
\subsection{Future work}

\section{Conclusion}
\label{sec:conclusion}
\todo{To be written after data collection.}

%% ---------------------------------------------------------------- disclosure
\section*{Use of AI}

Claude Code (Anthropic), running Claude Opus 5.5 with Claude Fable 5.1
subagents, was used in preparing the proposal that this paper grew from
to: critique the instructor's feedback and the original proposal; read the
full text of the MiPP paper and check our claims against it; check each
bibliographic entry against publisher records (Crossref, IEEE Xplore,
Springer, ACM, ICLR proceedings, arXiv) and produce corrected BibTeX; draft
the text of the proposal; and draw its mock-ups. For this checkpoint,
Claude Code (Claude Fable 5.1) ported the abstract, introduction and
related work into the review template, drafted the opening sentences of
the Method section from the released apparatus code and data, and built
the apparatus code described in Section~\ref{sec:method} with the authors.
The authors made every design decision, revise the text, and verify each
reference on the publisher page before submission; claims the tool could
not confirm are marked in the source and in the text.
% TODO authors: verify -- add any other AI tools used (for example, for
% grammar or LaTeX fixes) and what each did.

%% ---------------------------------------------------------------- references
\bibliographystyle{ACM-Reference-Format}
\bibliography{references}

\end{document}
